\documentclass[../Thesis]{subfiles}
\begin{document}
\chapter{Method} \label{chapter3}
\chaptermark{Method}
\section{Overview}
The methodology integrates dataset selection, preprocessing, model design, embedding extraction, similarity-matching, and evaluation. It is engineered to directly address the research objectives and limitations identified in Chapter 2 (\ref{chapter2}).

\section{Datasets}
This project uses open, and ethically cleared datasets, including the Free Music Archive and GTZAN \parencite{defferrard2017fma,tzanetakis2002genre}. Synthetic transformations will replicate real-world short-form edits.

\section{Audio Preprocessing}
All audio is resampled to a uniform rate, segmented into fixed windows, and converted to spectrograms using Librosa.

\section{Embedding Extraction}
Two pathways will be explored:
\begin{itemize}
\item Pre-trained embeddings (OpenL3, VGGish)
\item A contrastive model trained on augmented audio pairs
\end{itemize}
\subsection{Contrastive Learning}
A custom contrastive model is trained with distortion-focused augmentations: Tempo change, Pitch shifting, Cropping, Reverb, Noise injection and Mash-ups. This models the types of transformations common on TikTok, Instagram, and YouTube Shorts.

\section{Similarity Matching \& Evaluation}
Similarity matching in WaveID involves comparing audio embeddings to determine whether two clips originate from the same source track. Rather than relying on exact fingerprint matches, the system computes Cosine similarity and FAISS-based indexing \parencite{johnson2019faiss} that will be used to retrieve nearest-neighbour matches efficiently when working with large datasets. This approach supports scalable retrieval, allowing the system to compare millions of embedding vectors without sacrificing performance.

To evaluate the effectiveness of this similarity matching pipeline, the project adopts a comprehensive set of quantitative metrics. Precision, recall, and F1-score are used to measure how reliably the system identifies true matches while avoiding false positives. These metrics ensure a balanced understanding of accuracy, especially in scenarios where many clips may be only partially transformed or extremely short. 

Robustness testing forms a central part of the evaluation strategy. Each transformation explored in Chapter 2, pitch shifting, tempo changes, noise addition, cropping, and composite effects is applied systematically to assess how embedding quality degrades under increasingly extreme modifications. The system's ability to retrieve the correct reference track despite these distortions provides insight into the real-world viability of the model.

Scalability is also evaluated by gradually increasing the size of the dataset and observing how search latency and retrieval accuracy behave under realistic load. This is particularly important for audio identification systems, where operational environments (e.g., social-media platforms) must process millions of short clips daily. By combining a similarity search performance with distortion-based robustness testing, this integrated evaluation framework provides a holistic assessment of WaveID’s capability to operate effectively under practical constraints.

\section{Evaluation}
Evaluation includes:
\begin{itemize}
\item Precision, recall, F1-score
\item Robustness testing under pitch/tempo/noise/cropping conditions
\item Scalability testing on expanding dataset sizes
\end{itemize}

\section{Ethics}
This study uses only public, anonymised datasets and involves no human participants. GDPR does not apply.

\section{Research Methodology Justification}
Deep embeddings and contrastive learning were selected because they exhibit robustness to distortions common in short-form media. FMA and GTZAN are widely used in MIR research, allowing standardised comparisons. Evaluation metrics ensure replicability and quantitative benchmarking.

\section{Ethical, Social, and Legal Considerations}
This project complies with data protection, licensing, and institutional research ethics guidelines. It promotes responsible audio identification by restricting experiments to legally clear datasets.

\end{document}
